\begin{proposition}[4.5.19]
\EGAReaderTarget{ega4.c04.prop.4-5-19}{IV.2.4.5.19-fr}
Sean $X$ un $k$-preesquema, $Z$ una parte de $X$, $k'$ una extensión
algebraicamente cerrada de $k$, $X'=X\otimes_k k'$, $p:X'\to X$ la proyección
canónica. Se supone que $Z'=p^{-1}(Z)$ está contenida en \textup{una sola} componente
irreducible $X'_0$ de $X'$. Entonces se tiene $X'_0=p^{-1}(X_0)$, donde
$X_0=p(X'_0)$ es una componente irreducible de $X$ que contiene a $Z$, y además
$X_0$ es geométricamente irreducible. Si además $X'_0$ es la única componente
irreducible de $X'$ que \textup{corta} $Z'$, entonces $X_0$ es la única componente
irreducible de $X$ que corta $Z$.
\end{proposition}

Para demostrar que $X'_0=p^{-1}(X_0)$, vamos a aplicar el

\begin{lemma}[4.5.19.1]
\label{IV.2.4.5.19.1-fr}
Sean $T$ y $T'$ dos preesquemas, $p:T'\to T$ un morfismo~; consideremos el
preesquema $T''=T'\times_T T'$, y sean $p_1:T''\to T'$ y $p_2:T''\to T'$ las
proyecciones canónicas. Para que una parte $U'$ de $T'$ sea de la forma
$p^{-1}(U)$, donde $U\subset T$, es necesario y suficiente que se tenga
$p_1^{-1}(U')=p_2^{-1}(U')$.
\end{lemma}

\begin{proof}
El morfismo estructural $q:T''\to T$ es igual a $p\circ p_1$ y a
$p\circ p_2$ por definición, de modo que la relación es necesaria.
Recíprocamente, supongamos que se verifica~; sean $u'$ un punto de $U'$ y $t'$
un punto de $p^{-1}(p(u'))$~; como $p(t')=p(u')$, existe un punto $t''\in T''$
tal que $p_1(t'')=u'$ y $p_2(t'')=t'$ (\textbf I, 3.4.7)~; por tanto,
$t''\in p_1^{-1}(U')=p_2^{-1}(U')$ por hipótesis y, por consiguiente,
$t'=p_2(t'')\in U'$, lo que demuestra el lema.
\end{proof}

\begin{proof}
Demostrado este lema, para probar que $X'_0=p^{-1}(X_0)$, formemos el producto
$X''=X'\times_X X'$ (relativo al morfismo $p:X'\to X$)~; si $p_1$ y $p_2$ son
las proyecciones canónicas de $X''$ en $X'$, se trata, pues, de demostrar que
$p_1^{-1}(X'_0)=p_2^{-1}(X'_0)$. Ahora bien, si se pone
$S=\operatorname{Spec}(k)$ y $S'=\operatorname{Spec}(k')$, observemos que, si
$S''=S'\times_S S'=\operatorname{Spec}(k'\otimes_k k')$, también puede
escribirse $X''=X'\times_{S'}S''$. Si $\varphi'':X''\to S''$ es el morfismo
estructural, basta demostrar que, para todo $s''\in S''$, se tiene
$p_1^{-1}(X'_0)\cap\varphi''^{-1}(s'')
=p_2^{-1}(X'_0)\cap\varphi''^{-1}(s'')$.
Si se pone $T=\operatorname{Spec}(\boldsymbol{k}(s''))$,
$U=X''\times_{S''}T=\varphi''^{-1}(s'')$, y si $v:U\to X''$ es la proyección
canónica, se trata, pues, de demostrar que
$U_1=v^{-1}(p_1^{-1}(X'_0))$ y $U_2=v^{-1}(p_2^{-1}(X'_0))$, que son
componentes irreducibles de $U$ (4.4.5), son idénticas. Se tiene el diagrama
\begin{equation}
\tag{4.5.19.2}
\label{IV.2.4.5.19.2-fr}
\xymatrix@C=3.2em@R=3.0em{
X\ar[d]
  & X'\ar[l]_p\ar[d]
  & X''\ar@<.55ex>[l]^{p_2}\ar@<-.55ex>[l]_{p_1}
       \ar[d]^{\varphi''}
  & U\ar[l]_v\ar[d]
\\
S
  & S'\ar[l]
  & S''\ar@<.55ex>[l]\ar@<-.55ex>[l]
  & T\ar[l]
}
\end{equation}
Por construcción, tanto $U_1$ como $U_2$ contienen el conjunto
$w^{-1}(Z)$, donde
\[
w=p\circ p_1\circ v=p\circ p_2\circ v.
\]

\oldpage[IV-2]{67}

Ahora bien, existe un cuerpo $K$, extensión común de $k'$ y de $\boldsymbol{k}(s'')$, de
modo que, si se pone $P=\operatorname{Spec}(K)$, el diagrama
\[
\xymatrix@C=3.5em@R=3em{
S'\ar[d] & P\ar[l]\ar[d]\\
S & T\ar[l]
}
\]
sea conmutativo~; si se pone $Q=X\otimes_k K$, el diagrama
\[
\xymatrix@C=3.5em@R=3em{
X'\ar[d]_p & Q\ar[l]_q\ar[d]^r\\
X & U\ar[l]^w
}
\]
es, pues, también conmutativo, y las dos componentes irreducibles
$r^{-1}(U_1)$ y $r^{-1}(U_2)$ de $Q$ (4.4.5) contienen, por consiguiente,
$q^{-1}(Z')$. Ahora bien, en virtud de (4.4.5), las imágenes
$q(r^{-1}(U_1))$ y $q(r^{-1}(U_2))$ son componentes irreducibles de $X'$ que
contienen a $Z'$~; por tanto, ambas son iguales a $X'_0$, luego
$r^{-1}(U_1)=r^{-1}(U_2)$ en virtud de (4.4.5), y finalmente $U_1=U_2$.

Como $p$ es sobreyectivo y $X'_0$ domina una componente irreducible de $X$,
$X_0=p(X'_0)$ es una componente irreducible de $X$ que contiene a $Z$, y como
$p^{-1}(X_0)$ es una componente irreducible de $X'$, $X_0$ es geométricamente
irreducible.

Para demostrar la última afirmación, sea $X_1$ una segunda componente
irreducible de $X$ que corte $Z$. Tomando $z\in Z\cap X_1$, se ve que existe
una componente irreducible de $X'$ que corta $Z'$ y domina $X_1$ (2.3.5)~;
pero, como $X'_0$ es por hipótesis la única componente irreducible de $X'$ que
corta $Z'$, se tiene necesariamente $X_1=X_0$.
\end{proof}

\begin{remark}[4.5.20]
\label{IV.2.4.5.20-fr}
El hecho de que $Z'$ esté contenida en una sola componente irreducible $X'_0$
de $X'$ \emph{no implica} que $X_0=p(X'_0)$ sea la única componente irreducible
de $X$ que contiene a $Z$. Tomemos, en efecto, $k=\mathbf R$, $k'=\mathbf C$, y
como $X$ el $k$-esquema obtenido por pegado\footnote{La técnica del
«~pegado~» de preesquemas se desarrollará en detalle en el cap.~V~; para el
caso considerado aquí (donde se trata de curvas algebraicas), puede consultarse
[38], pp.~68--71.} en un punto (no genérico) de
$X_1=\operatorname{Spec}(\mathbf R[S,T]/(S^2+T^2+1))$ y de
$X_2=\operatorname{Spec}(\mathbf C[T])$ (los anillos locales de estas dos
curvas en cada uno de sus puntos cerrados tienen el mismo cuerpo residual
isomorfo a $\mathbf C$)~; $X_1$ y $X_2$ se identifican con las dos componentes
irreducibles de $X$~; si $x$ es su punto común, $p^{-1}(x)$ está formado por
dos puntos distintos $y'$ y $z'$, $p^{-1}(X_1)$ es irreducible, y
$p^{-1}(X_2)$ es suma de dos componentes irreducibles $Y'$ y $Z'$ de $X'$
tales que $y'\in Y'$ y $z'\in Z'$, de modo que $p^{-1}(x)$ solo está contenido
en una única componente irreducible de $X'$.
\end{remark}

\begin{proposition}[4.5.21]
\label{IV.2.4.5.21-fr}
Sea $k$ un cuerpo separablemente cerrado, y sea $k'$ una clausura algebraica de
$k$. Para todo $k$-preesquema $X$, la proyección canónica
$X\otimes_k k'\to X$ es un homeomorfismo universal.

\oldpage[IV-2]{68}

En particular, las componentes irreducibles \textup{(resp.} conexas\textup{)} de $X$ son
geométricamente irreducibles \textup{(resp.} geométricamente conexas\textup{)}.
\end{proposition}

\begin{proof}
Por definición, $k'$ es una extensión radicial de $k$~; el morfismo
$\operatorname{Spec}(k')\to\operatorname{Spec}(k)$ es entero, sobreyectivo y
radicial, por lo que la primera afirmación se deduce de (2.4.5, (i))~; la
segunda se deduce de ella, en vista de (4.5.1).
\end{proof}

